\documentclass{article}
\usepackage{graphicx} % Required for inserting images
\usepackage{optidef} % Helps for optimization formulation
\usepackage{amssymb} % ADDED: for \mathcal
\usepackage{booktabs} % ADDED: for the scenario table

\title{Cost-Optimal Data Center Self-Supply Under Curtailment}
\author{Charles Loriot, Conor McManamy}
\date{November 10, 2026}

\begin{document}

\maketitle

\section{Background}

This model captures a stochastic two-stage capacity expansion exercise to illustrate the cost-optimal backup power supply that data centers need to procure to tolerate a range of expected curtailment events.

\section{Sets and Indices}

\begin{itemize}
    \item Generator build candidates: ${g \in G }$, with subsets such as $G^{BESS}$, $G^{DIES}$, $G^{VPP}$, and $G^{PV}$.
    \item Built generators: $G_{built} \subset G$.
    \item Plausible curtailment scenarios that could occur during a planning horizon $H$, each with varying duration: ${\omega \in {\Omega}}$
    \begin{itemize}
        \item $\Omega = H^{0} \times \mathcal{L} \times \mathcal{K}$: each scenario combines a start hour, a duration and a depth, giving $3 \times 7 \times 4 = 84$ scenarios (see Section~3.1) % ADDED
        \item Start hours: $h^{0}_\omega \in H^{0} = \{4, 5, 6 \text{ p.m.}\}$ % ADDED
        \item Durations: $L_\omega \in \mathcal{L} = \{3, 4, 5, 6, 7, 8, 9\}$ hours % ADDED
        \item Depths: $\kappa_\omega \in \mathcal{K} = \{0.25, 0.50, 0.75, 1.0\}$ % ADDED
    \end{itemize}
    \item Time steps within any curtailment scenario $\omega$: $t \in \{1,2, ..., {T}\}$, with $T = L_\omega$ the duration of scenario $\omega$ % ADDED: link T to duration
\end{itemize}


\section{Parameters}

\subsection{Curtailment scenarios} % ADDED (whole subsection)

Scenarios aim to be realistic for an ERCOT site. Each scenario $\omega$ is one
curtailment event, defined by its start hour, its duration and its depth.

\begin{center}
\small
\begin{tabular}{ll}
\toprule
Dimension & Values \\
\midrule
Start hour $h^{0}_\omega$ & 4, 5, 6 p.m. \\
Duration $L_\omega$ & 3 to 9 h, in 1 h steps \\
Depth $\kappa_\omega$ & 0.25, 0.50, 0.75, 1.0 \\
\bottomrule
\end{tabular}
\end{center}

\begin{itemize}
    \item \textbf{Weights.} All 84 combinations of start hour, duration and depth are equally likely: $w_\omega = 1/|\Omega| = 1/84$.
    \item \textbf{Events per year.} $N = 19$ curtailment events per year (Aurora 2026, North Hub), with $N \in \{6, 38\}$ as sensitivities.
    \item \textbf{Depth.} Each scenario has one depth, constant over the event: $\kappa_{\omega,t} = \kappa_\omega$ for all $t$. 
\end{itemize}

\noindent Availability $\alpha_g$ of each resource during a scenario:
\begin{itemize}
    \item Diesel ($G^{DIES}$): runs as long as there is fuel ($E_g$).
    \item Storage, including long-duration ($G^{BESS}$): runs as long as there is charge ($E_g$).
    \item Solar ($G^{PV}$): runs on a representative site-agnostic solar output profile.
    \item VPP ($G^{VPP}$): participation rate set by the time of day, then degrading linearly with the duration of the event, $\alpha^{VPP}_t = \alpha^{VPP}_{hr} - \delta(t)$.
    \item Availability does not change across scenarios beyond these hourly profiles.
\end{itemize}

\section{Objective Function}
% ADDED (whole section)

\begin{equation}
\begin{split}
\min_{P, E, \theta, b^{E}, \sigma, u} \quad
& \sum_{g \in G} \left( C^{P}_{g} P_g + C^{E}_{g} E_g \right) \\
& + N \sum_{\omega \in \Omega} w_\omega \sum_{t=1}^{T} \Big[ \sum_{g \in G} \left( \pi^{MC}_g + \pi^{c} e_g \right) \theta_{g,\omega,t} + \pi^{UE} u_{\omega,t} \Big]
\end{split}
\end{equation}

\section{Constraints}
% ADDED (whole section, from the constraints in the Model block). G^E = G^DIES and G^BESS (candidates with an energy limit).

\noindent Here $G^{E} = G^{DIES} \cup G^{BESS}$ are the candidates with an energy limit (fuel or charge), and $t$ runs over the hours of each scenario.

\begin{align}
& (\kappa_{\omega,t} - \sigma_{\omega,t}) D_{\omega,t} = \sum_{g \in G} \theta_{g,\omega,t} + u_{\omega,t} && \forall \omega, t \\
& 0 \le P_g \le P^{max}_g && \forall g \in G \\
& 0 \le E_g \le E^{max}_g && \forall g \in G^{E} \\
& 0 \le \theta_{g,\omega,t} \le \alpha_{g,\omega,t} P_g && \forall g \in G, \omega, t \\
& b^{E}_{g,\omega,1} = E_g && \forall g \in G^{E}, \omega \\
& b^{E}_{g,\omega,t+1} = b^{E}_{g,\omega,t} - \theta_{g,\omega,t} / \eta_g && \forall g \in G^{E}, \omega, t \\
& 0 \le b^{E}_{g,\omega,t} \le E_g && \forall g \in G^{E}, \omega, t \\
& 0 \le \sigma_{\omega,t} \le \sigma^{max} && \forall \omega, t \\
& u_{\omega,t} \ge 0 && \forall \omega, t
\end{align}

\begin{enumerate}
    \item Power balance: in every hour of an event, backup output plus unserved load covers the curtailed load that is not shifted.
    \item Power build limit: the MW built or contracted stays within what can be sited or procured.
    \item Energy build limit: the fuel tank or storage size (MWh) stays within what can be built.
    \item Availability: output cannot exceed the available share of capacity (solar decline, VPP participation, outages).
    \item Starting level: fuel and storage start every event full.
    \item Energy tracking: fuel or charge drops each hour by the output divided by its efficiency.
    \item Energy bounds: fuel or charge stays between empty and the size built.
    \item Load shift limit: the data center can shift at most a set share of its load.
    \item Unserved load cannot be negative.
\end{enumerate}

\section{Problem class, solver, and other details}

\textbf{Problem class.} Two-stage stochastic linear program (LP). The first stage determines investments in capacity and energy resources and incurs annualized investment and fixed O\&M costs. The first stage determines the build of each backup option ($P_g$, $E_g$). The second stage averages cost for operation ($\theta$, $b^{E}$, $\sigma$, $u$) throughout each curtailment event scenario. All scenarios are solved together in one model, and since every cost and constraint is linear, the problem is an LP and solves exactly. 

\textbf{Solver.} Gurobi, a commercial mathematical optimization tool. 

\section{Exclusions}

The model makes the following assumptions:

\begin{enumerate}
    \item fuel and stored state of charge are always at peak at the beginning of each curtailment scenario.
    \item events are small enough such that load shifting can defer workload outside of the curtailment event without violating quality-of-service.
    \item resources have negligible start-up costs on the hourly timescale.
    \item events are sufficiently rare so as to allow a full recovery between them. In other words, events are standalone, not back-to-back. 
    \item capacity expansion decisions are continuous. Discreteness would force a MILP formulation.
    \item no grid conditions are considered, beyond those proxied for by each curtailment event.
\end{enumerate}



\section{Model}

\begin{mini*}|s|
{P, E, \theta, b^{E}, \sigma, u}
{\sum_{g \in G} \left( C^{P}_{g} P_g + C^{E}_{g} E_g \right)}
{}{}
\breakObjective{+ N \sum_{\omega \in \Omega} w_\omega \sum_{t=1}^{T} \Big[ \sum_{g \in G} \left( \pi^{MC}_g + \pi^{c} e_g \right) \theta_{g,\omega,t} + \pi^{UE} u_{\omega,t} \Big]}
\addConstraint{(\kappa_{\omega,t} - \sigma_{\omega,t}) D_{\omega,t}}{= \sum_{g \in G} \theta_{g,\omega,t} + u_{\omega,t}}{\quad \forall \omega, t}
\addConstraint{0 \le P_g}{\le P^{max}_g}{\quad \forall g \in G}
\addConstraint{0 \le E_g}{\le E^{max}_g}{\quad \forall g \in G^{E}}
\addConstraint{0 \le \theta_{g,\omega,t}}{\le \alpha_{g,\omega,t} P_g}{\quad \forall g \in G, \omega, t}
\addConstraint{b^{E}_{g,\omega,1}}{= E_g}{\quad \forall g \in G^{E}, \omega}
\addConstraint{b^{E}_{g,\omega,t+1}}{= b^{E}_{g,\omega,t} - \theta_{g,\omega,t} / \eta_g}{\quad \forall g \in G^{E}, \omega, t}
\addConstraint{0 \le b^{E}_{g,\omega,t}}{\le E_g}{\quad \forall g \in G^{E}, \omega, t}
\addConstraint{0 \le \sigma_{\omega,t}}{\le \sigma^{max}}{\quad \forall \omega, t}
\addConstraint{u_{\omega,t}}{\ge 0}{\quad \forall \omega, t}
\end{mini*}

\end{document}
